\chapter{HPC Execution and Parallelization Boundary}
\label{ch:hpc}

\section{HPC cluster execution environment}
Production finite element simulations and remeshing workflows are executed on the high-performance computing cluster at TU Bergakademie Freiberg. The verified software environment comprises:
\begin{itemize}
  \item \textbf{Finite Element Solver:} \Abaqus/Standard 2023 (hotfix-level 2023.HF4);
  \item \textbf{Compiler Toolchain:} Intel oneAPI Fortran Compiler (ifx 2024.2.0 / ifort compatibility);
  \item \textbf{Workload Manager:} PBS Professional batch queuing system;
  \item \textbf{Compute Nodes:} Dedicated compute nodes (e.g., \texttt{mnode097}, \texttt{mnode099}) equipped with dual-socket Intel Xeon processors and 128--256\,GB RAM;
  \item \textbf{Standard Allocation Directives:} \texttt{nodes=1:ppn=1}, 16--32\,GB requested RAM, queue \texttt{normal\_imfdfkmq}.
\end{itemize}

Operational job identifiers, execution checksums, walltimes, and resource allocations are cataloged in Appendix~\ref{app:repro}.

\section{Serial execution boundary and state synchronization}
A key technical finding of this project concerns the parallelization safety of the coupled phase-field UEL implementation.

In the Moln\'ar--Gravouil staggered architecture \citep{Molnar2017}, the phase-field UEL and mechanical UEL layers exchange trial strains, driving energy $\mathcal{H}$, and phase damage $d$ through shared Fortran global arrays in common storage:
\begin{verbatim}
COMMON /CB_STATE_TRANS/ SV_PHASE_TRIAL(MAXELEM), SV_HIST_TRIAL(MAXELEM, 4), ...
\end{verbatim}
During an Abaqus solver iteration, element stiffness and residual evaluations are executed in an element-by-element loop. In a shared-memory multithreaded execution (e.g., \texttt{cpus=4} or OpenMP threading), multiple worker threads execute the UEL subroutine concurrently for different element ranges. Because the common block storage lacks explicit thread-level mutual exclusion locks or memory barriers, concurrent read and write operations on shared element indices create race conditions and non-deterministic memory corruption.

In distributed-memory MPI execution (e.g., domain decomposition across ranks), each MPI rank runs in a private address space; common blocks are replicated per process rather than shared globally, causing ranks to lose access to elements residing on other partitions.

Consequently, all production calculations in this research are governed by a strict \textbf{single-core serial execution boundary} (\texttt{cpus=1}). A prospective multithreaded or MPI implementation requires refactoring state storage from global common blocks to Abaqus-managed user state tensors (\texttt{STATEV}) or explicit MPI communication buffers.

\section{Notification path verification}
To ensure reliable monitoring of long-running cluster jobs (which require up to 35 hours of walltime), automated notification traps were integrated into the PBS job submission scripts:
\begin{itemize}
  \item \textbf{Telegram Bot Webhook Trap:} Configured to transmit job start, step transitions, walltime milestones, and terminal exit codes ($0$ vs error) directly to the engineering team. This notification channel has been verified as \textbf{ACTIVE AND RELIABLE}.
  \item \textbf{Direct SMTP Email Notification (\texttt{\#PBS -m abe}):} PBS mail delivery exhibited an internal relay failure on the cluster network and is recorded as a \textbf{KNOWN INFRASTRUCTURE DEFECT}.
\end{itemize}
In accordance with project governance rules, a running scientific job is never cancelled or invalidated due to email delivery failures when the calculation itself executes normally.
